\section{Discussion and Limitations}
\label{sec:discussion}

\subsection{When Does Deep Learning Help, and When Does It Not?}
We compare classical machine learning on engineered descriptors with deep learning on molecular graphs and report what the results do and do not support.

\textbf{1. Graph structure versus flat inputs: weak evidence.} The H1 hybrid (\texttt{mpnn+desc}) has a higher observed mean scaffold PR~AUC than the FP-MLP control ($0.3553$ vs.\ $0.3429$, $+0.0124$) and a higher PR~AUC on 9 of 12 endpoints, largest on \texttt{NR-Aromatase} ($+0.0980$), but the paired Wilcoxon test gives $p = 0.1514$, so the difference is not significant. The two models also differ in their inputs (the FP-MLP sees Morgan bits plus descriptors, H1 sees the graph plus descriptors). The graph-only models score below the FP-MLP in mean scaffold PR~AUC (GCN $0.3397$, MPNN $0.3069$).

\textbf{2. The GCN beat the plain MPNN.} The GCN scores higher than the MPNN on scaffold ($0.3397$ vs.\ $0.3069$ PR~AUC), with a higher PR~AUC on 11 of 12 endpoints (mean difference $+0.0328$, $p = 0.0015$), and has a smaller seed standard deviation ($0.0027$ vs.\ $0.0180$). The more expressive edge-network MPNN therefore did not improve on the simpler baseline here; we did not investigate why.

\textbf{3. Classical tree ensembles remain best on scaffold.} Random Forest ($0.4035$ PR~AUC) and XGBoost ($0.4010$) on combined fingerprints and descriptors score above every deep model and above both H2 variants on the scaffold split. This is in line with the comparison study of Jiang et al.~\cite{jiang2021could}, who report that descriptor-based models are often competitive with graph neural networks. On the random split the primary H2 has the highest mean PR~AUC ($0.5727$ vs.\ $0.5647$ for XGBoost), without a significant difference ($9/3$ endpoint wins, $p = 0.0923$).

\textbf{4. Descriptors as global information.} The frozen MPNN has $628{,}300$ trainable parameters (counted from the code with the frozen configuration) and is trained on about $6{,}200$ molecules. Descriptors such as lipophilicity or polar surface area encode whole-molecule properties directly, which may help a model trained on little data; we did not test this explanation. What the shuffled-descriptor control shows is narrower: with the same architecture, permuting descriptor rows across molecules lowers H1 from $0.3553$ to $0.2547$ PR~AUC, so its benefit depends on correct alignment between molecules and descriptor values. The control does not identify which descriptors matter or why.

\textbf{5. Ensembling.} The primary H2 (XGBoost plus FP-MLP, partner selected on validation PR~AUC) scores $0.3945$ on scaffold. It beats its deep components (against the FP-MLP on 12 of 12 endpoints, $p = 0.0005$; against H1 on 11 of 12, $p = 0.0093$) but not XGBoost alone: it wins on 4 endpoints and loses on 8 (mean difference $-0.0065$, $p = 0.5186$). With equal-weight rank averaging, fusing a deep model with XGBoost therefore did not improve on XGBoost alone on the scaffold split. We did not test other fusion weights, and we make no claim about error complementarity.

\subsection{Methodological Limitations}
Our protocol freezes the configurations before test results are committed, but several limitations remain:

\begin{itemize}
    \item \textbf{Small Dataset Size and Single Benchmark:} Our experiments are conducted exclusively on the Tox21 benchmark ($N = 7,709$). In massive screening libraries ($N > 100,000$) or under self-supervised pre-training regimes, deep graph architectures typically exhibit stronger scaling advantages that cannot be realized on small datasets.
    \item \textbf{Endpoint Sample Size and Statistical Power:} With $K = 12$ endpoints, non-parametric statistical tests (such as the Wilcoxon signed-rank test) possess inherently limited statistical power. Bootstrap intervals on the seed-0 predictions exclude zero on only 1 of 12 endpoints for H1 minus FP-MLP, so endpoint-level evidence is weaker than the pooled win counts. More assays would give tighter bounds.
    \item \textbf{Hyperparameter Budget Constraints:} Hyperparameter search was limited to at most 8 configurations per model, each evaluated once with seed 0 on the scaffold validation split, and the scaffold choice is reused unchanged on the random split. H1 has no search of its own: it reuses the frozen MPNN configuration. A larger search might change the ranking of the deep models.
    \item \textbf{Ablations run on the test split.} The ablations were evaluated on the test split, so they are explanatory and were not used to choose any configuration. The command lines of the ablation runs are not recorded in the repository and were reconstructed from the result files. In the shuffled-label control (Section~\ref{sec:malhar-architecture-ablations}) the seed-2 test ROC~AUC ($0.5632$) is above the predeclared $0.55$ threshold; the cause has not been investigated.
    \item \textbf{2D Topological Graphs vs.\ 3D Conformers:} Our graph formulation is restricted to 2D topological connectivity and atom/bond features. Toxicity bioassays frequently depend on 3D stereochemical fit into biological receptor binding pockets (e.g., nuclear receptor ligand binding domains). Incorporating 3D conformer ensembles or quantum chemical surface potentials represents an open technical frontier.
\end{itemize}
